% !TeX program = xelatex
% Báo cáo kết thúc môn (mở rộng 30-50 trang) — từ dự án hiện tại

\documentclass[12pt,a4paper,oneside]{report}

% =========================
% 0) Thiết lập chung
% =========================
\usepackage[a4paper,margin=2.5cm]{geometry}
\usepackage{setspace}
\onehalfspacing
\usepackage{indentfirst}
\setlength{\parskip}{6pt}
\setlength{\parindent}{1.2cm}

% =========================
% 1) Tiếng Việt + Font (XeLaTeX)
% =========================
\usepackage{fontspec}
\usepackage{polyglossia}
\setdefaultlanguage{vietnamese}

\IfFontExistsTF{Times New Roman}{
  \setmainfont{Times New Roman}
}{
  \setmainfont{TeX Gyre Termes}
}

\IfFontExistsTF{Arial}{
  \setsansfont{Arial}
}{
  \setsansfont{TeX Gyre Heros}
}

\IfFontExistsTF{JetBrains Mono}{
  \setmonofont{JetBrains Mono}
}{
  \IfFontExistsTF{Consolas}{
    \setmonofont{Consolas}
  }{
    \setmonofont{Latin Modern Mono}
  }
}

% =========================
% 2) Gói tiện ích
% =========================
\usepackage{graphicx}
\usepackage{float}
\usepackage{booktabs}
\usepackage{array}
\usepackage{multirow}
\usepackage{longtable}
\usepackage{enumitem}
\usepackage{caption}
\usepackage{subcaption}
\usepackage{xcolor}
\usepackage{hyperref}
\usepackage{url}
\usepackage{fancyhdr}
\usepackage{listings}
\usepackage{amsmath}
\usepackage{amssymb}
\usepackage{microtype}

\hypersetup{
  colorlinks=true,
  linkcolor=black,
  urlcolor=blue,
  citecolor=black
}

\pagestyle{fancy}
\fancyhf{}
\lhead{Báo cáo kết thúc môn}
\rhead{Cổng thông tin Sa Đéc}
\cfoot{\thepage}

\lstdefinestyle{codestyle}{
  basicstyle=\ttfamily\small,
  keywordstyle=\color{blue!70!black},
  commentstyle=\color{gray!70!black},
  stringstyle=\color{teal!60!black},
  breaklines=true,
  frame=single,
  rulecolor=\color{black!20},
  tabsize=2,
  showstringspaces=false,
  language=csh
}
\lstset{style=codestyle}

% =========================
% 3) Thông tin dự án
% =========================
\newcommand{\ProjectName}{Cổng thông tin quảng bá Sa Đéc}
\newcommand{\FrontendUrl}{https://ashy-field-09c183d00.2.azurestaticapps.net}
\newcommand{\BackendUrl}{https://sadec-backend.azurewebsites.net}
\newcommand{\SwaggerUrl}{https://sadec-backend.azurewebsites.net/swagger}

% Macro placeholder
\newcommand{\ImagePlaceholder}[3]{%
  \begin{figure}[H]
    \centering
    \fbox{\parbox{0.92\textwidth}{\vspace{18pt}\centering\textbf{[CHÈN ẢNH]} \\ \textbf{#1}\\\vspace{6pt}\textit{Tên file:} \texttt{#2}\\\vspace{6pt}\textit{Nội dung:} #3\vspace{18pt}}}
  \end{figure}
}

\begin{document}

% =========================
% TRANG BÌA (TỰ THIẾT KẾ)
% =========================
% % Trang bìa bạn tự làm theo mẫu trường.
% Gợi ý: dùng \thispagestyle{empty} và canh giữa.
% File này đang để trống theo yêu cầu.

\begin{titlepage}
\thispagestyle{empty}
\centering
\vspace*{2cm}
{\Large \textbf{(TRANG BÌA TỰ THIẾT KẾ)}}\\
\vspace{1cm}
{\large \textbf{\ProjectName}}\\
\vfill
{\large \today}\\
\end{titlepage}
\clearpage


% =========================
% MỤC LỤC
% =========================
\pagenumbering{roman}
\tableofcontents
\listoffigures
\listoftables
\newpage

% =========================
% CHƯƠNG 1: GIỚI THIỆU
% =========================
\pagenumbering{arabic}

\chapter{Giới thiệu}

\section{Bối cảnh và lý do chọn đề tài}

Trong bối cảnh chuyển đổi số tại các địa phương, nhu cầu cung cấp thông tin công khai và quảng bá du lịch ngày càng tăng.
Thành phố Sa Đéc, một thành phố truyền thống giàu tiềm năng về du lịch và nền văn hóa, cần có một cổng thông tin định tập trung
để quản lý các thông tin: tin tức địa phương, điểm đến du lịch, ẩm thực đặc sắc, và các sự kiện quan trọng.

Một hệ thống web hiện đại không chỉ giúp cung cấp thông tin kịp thời mà còn tạo cơ hội thu hút du khách, phát triển du lịch
và nâng cao hình ảnh của địa phương trên các nền tảng kỹ thuật số.

\section{Mục tiêu đề tài}

Đề tài này hướng tới các mục tiêu sau:

\begin{enumerate}[itemsep=4pt]
  \item Xây dựng website công khai (public) với giao diện thân thiện người dùng, hiển thị:
  \begin{itemize}
    \item Tin tức địa phương (news).
    \item Điểm đến du lịch (destinations).
    \item Ẩm thực đặc sắc (cuisines).
    \item Bản đồ hành chính và thông tin liên hệ.
    \item Feedback từ khách hàng.
  \end{itemize}
  
  \item Xây dựng hệ thống quản trị (admin backend) cho phép:
  \begin{itemize}
    \item Biên tập viên tạo, cập nhật, xoá nội dung.
    \item Quản lý trạng thái nội dung (draft, published, archived).
    \item Upload và quản lý hình ảnh/media.
    \item Duyệt bình luận và phản hồi.
  \end{itemize}
  
  \item Đảm bảo bảo mật:
  \begin{itemize}
    \item Xác thực người dùng bằng JWT (JSON Web Token).
    \item Phân quyền theo vai trò (Admin, Editor, Viewer).
    \item Bảo vệ API bằng CORS và authorization.
  \end{itemize}
  
  \item Triển khai trên nền tảng cloud (Azure):
  \begin{itemize}
    \item Frontend trên Azure Static Web Apps.
    \item Backend trên Azure App Service.
    \item Dữ liệu lưu trên Azure SQL Database.
  \end{itemize}
  
  \item Đảm bảo tính ổn định và tính sẵn sàng:
  \begin{itemize}
    \item Dữ liệu bền vững (persistent) trên Azure SQL.
    \item Tự động migration schema.
    \item Recovery mechanism cho các lỗi tạm thời.
  \end{itemize}
\end{enumerate}

\section{Phạm vi thực hiện}

Toàn bộ dự án bao gồm:

\subsection{Frontend}
\begin{itemize}[itemsep=2pt]
  \item Trang chủ public với tính năng search.
  \item Trang chi tiết nội dung (news, destinations, cuisines).
  \item Trang login/đăng xuất.
  \item Dashboard admin với các bảng quản lý CRUD.
  \item Trang cài đặt hệ thống (logo, cấu hình).
  \item Trang quản lý người dùng.
  \item Trang quên/đặt lại mật khẩu.
\end{itemize}

\subsection{Backend}
\begin{itemize}[itemsep=2pt]
  \item API public: lấy danh sách, chi tiết nội dung, cài đặt.
  \item API admin: CRUD nội dung, quản lý media, quản lý người dùng.
  \item API xác thực: login, refresh token, forgot/reset password.
  \item API quản lý cài đặt và log hành động admin.
\end{itemize}

\subsection{Cơ sở dữ liệu}
\begin{itemize}[itemsep=2pt]
  \item Các bảng entity chính: Users, News, Destinations, Cuisines, Comments, Settings, MediaFiles.
  \item Audit log cho các hành động quan trọng.
\end{itemize}

\section{Đối tượng sử dụng}

\begin{itemize}[itemsep=3pt]
  \item \textbf{Người dùng công khai (Public User):} Truy cập website để xem thông tin, tìm kiếm địa điểm, xem tin tức.
  \item \textbf{Biên tập viên (Editor):} Tạo, chỉnh sửa nội dung (news, destinations); không được quản lý người dùng.
  \item \textbf{Quản trị viên (Admin):} Toàn quyền: CRUD nội dung, quản lý người dùng, cài đặt hệ thống, xem log.
\end{itemize}

\section{Môi trường triển khai}

Hệ thống được triển khai trên các dịch vụ Azure:

\begin{center}
\begin{tabular}{|l|l|l|}
\hline
\textbf{Thành phần} & \textbf{Dịch vụ} & \textbf{URL/Endpoint} \\
\hline
Frontend & Azure Static Web Apps & \url{\FrontendUrl} \\
\hline
Backend API & Azure App Service & \url{\BackendUrl} \\
\hline
Swagger UI & App Service & \url{\SwaggerUrl} \\
\hline
Cơ sở dữ liệu & Azure SQL Database & sqlserver.database.windows.net \\
\hline
\end{tabular}
\end{center}

\ImagePlaceholder{Giao diện trang chủ}{ui_homepage.png}{Chụp trang chủ sau deploy: banner, tin tức nổi bật, địa điểm nổi bật, bản đồ}

% =========================
% CHƯƠNG 2: CƠ SỞ LÝ THUYẾT
% =========================

\chapter{Cơ sở lý thuyết}

\section{Kiến trúc ứng dụng web hiện đại}

\subsection{Mô hình Single Page Application (SPA)}

Ứng dụng SPA là kiểu ứng dụng web tải toàn bộ code và resource một lần, sau đó điều hướng bằng JavaScript mà không cần tải lại trang.
Ưu điểm:
\begin{itemize}[itemsep=2pt]
  \item Trải nghiệm người dùng mượt mà (smooth UX).
  \item Giảm tải cho server (chỉ phục vụ API).
  \item Cho phép offline hoạt động (với service worker).
\end{itemize}

Thách thức:
\begin{itemize}[itemsep=2pt]
  \item Reload trực tiếp route (ví dụ: \texttt{/admin/users}) có thể gặp lỗi 404 nếu server không cấu hình rewrite.
  \item Cần cấu hình CORS giữa frontend và backend nếu khác domain.
\end{itemize}

\subsection{API RESTful}

REST (Representational State Transfer) sử dụng HTTP method để biểu diễn các thao tác:

\begin{center}
\begin{tabular}{|l|l|l|l|}
\hline
\textbf{Method} & \textbf{Ý nghĩa} & \textbf{Endpoint ví dụ} & \textbf{Trả về} \\
\hline
GET & Lấy dữ liệu & \texttt{/api/news/1} & 200 OK + data \\
\hline
POST & Tạo mới & \texttt{/api/admin/news} & 201 Created \\
\hline
PUT & Cập nhật toàn bộ & \texttt{/api/admin/news/1} & 200 OK \\
\hline
PATCH & Cập nhật một phần & \texttt{/api/admin/news/1} & 200 OK \\
\hline
DELETE & Xoá & \texttt{/api/admin/news/1} & 204 No Content \\
\hline
\end{tabular}
\end{center}

\section{Xác thực và phân quyền (Authentication \& Authorization)}

\subsection{JWT (JSON Web Token)}

JWT là chuỗi mã hóa gồm 3 phần (header.payload.signature), cho phép server xác minh danh tính của client mà không cần lưu trạng thái session.

Cấu trúc JWT:
\begin{lstlisting}[language=json,caption={Ví dụ JWT Payload}]
{
  "sub": "user_id",
  "email": "editor@sadec.local",
  "roles": ["Editor"],
  "exp": 1712345600,
  "iat": 1712342000
}
\end{lstlisting}

Luồng xác thực:
\begin{enumerate}[itemsep=2pt]
  \item Client gửi username + password tới \texttt{/api/auth/login}.
  \item Server trả access token (JWT) + refresh token.
  \item Client gửi kèm header: \texttt{Authorization: Bearer <access\_token>}.
  \item Khi access token hết hạn, client dùng refresh token để lấy access token mới.
\end{enumerate}

\subsection{Role-Based Access Control (RBAC)}

Hệ thống sử dụng 3 vai trò chính:

\begin{center}
\begin{tabular}{|l|p{4cm}|p{4.5cm}|}
\hline
\textbf{Vai trò} & \textbf{Quyền} & \textbf{API được phép} \\
\hline
Admin & Toàn quyền & Tất cả endpoint \\
\hline
Editor & CRUD nội dung & \texttt{/api/admin/news/*}, \texttt{/api/admin/destinations/*} \\
\hline
Viewer & Chỉ đọc & \texttt{/api/public/*} \\
\hline
\end{tabular}
\end{center}

\section{Công nghệ Stack}

\subsection{Frontend Stack}
\begin{itemize}[itemsep=2pt]
  \item \textbf{React 18:} Thư viện UI component-based.
  \item \textbf{TypeScript:} Tăng type-safety và đơn giản hóa debugging.
  \item \textbf{Vite:} Build tool nhanh, HMR (Hot Module Replacement) tức thời.
  \item \textbf{React Router v6:} Định tuyến SPA.
  \item \textbf{Axios:} HTTP client với interceptor.
  \item \textbf{Tailwind CSS:} Styling framework utility-first.
\end{itemize}

\subsection{Backend Stack}
\begin{itemize}[itemsep=2pt]
  \item \textbf{ASP.NET Core 8:} Framework web API .NET.
  \item \textbf{ASP.NET Identity:} Quản lý user, role, password.
  \item \textbf{Entity Framework Core:} ORM cho SQL.
  \item \textbf{SQL Server / Azure SQL:} Cơ sở dữ liệu quan hệ.
  \item \textbf{Swagger:} Tài liệu API interactiv.
\end{itemize}

\subsection{DevOps \& Cloud}
\begin{itemize}[itemsep=2pt]
  \item \textbf{Azure Static Web Apps:} Triển khai frontend với CI/CD tự động.
  \item \textbf{Azure App Service:} Chạy backend ASP.NET.
  \item \textbf{Azure SQL Database:} Cơ sở dữ liệu cloud.
  \item \textbf{GitHub Actions:} Automation build và deploy.
\end{itemize}

\section{Entity-Relationship Model (ERD)}

\subsection{Các bảng chính}

\begin{itemize}[itemsep=3pt]
  \item \textbf{AspNetUsers} (từ ASP.NET Identity):
  \begin{itemize}
    \item Lưu user: Id, UserName, Email, PasswordHash, CreatedAt.
  \end{itemize}
  
  \item \textbf{AspNetRoles}: Admin, Editor, Viewer.
  
  \item \textbf{News} (Tin tức):
  \begin{itemize}
    \item Title, Slug, Content, Excerpt, FeaturedImageUrl.
    \item Status (Draft, Published, Archived).
    \item CreatedBy, CreatedAt, UpdatedAt.
  \end{itemize}
  
  \item \textbf{Destinations} (Địa điểm):
  \begin{itemize}
    \item Name, Slug, Description, ImageUrl.
    \item Coordinates (latitude, longitude).
    \item Status, CreatedAt.
  \end{itemize}
  
  \item \textbf{Cuisines} (Ẩm thực):
  \begin{itemize}
    \item Name, Slug, Description, ImageUrl.
    \item Status, CreatedAt.
  \end{itemize}
  
  \item \textbf{Comments} (Bình luận):
  \begin{itemize}
    \item Content, Status (Pending, Approved, Rejected).
    \item EntityType (News, Destination), EntityId.
  \end{itemize}
  
  \item \textbf{Settings} (Cài đặt hệ thống):
  \begin{itemize}
    \item Key (ví dụ: site.logo, site.title, site.administrativeMapImageUrl).
    \item Value (JSON).
  \end{itemize}
  
  \item \textbf{AuditLog} (Log hành động):
  \begin{itemize}
    \item Action, UserId, EntityType, EntityId, Timestamp.
  \end{itemize}
\end{itemize}

\ImagePlaceholder{ERD (Entity Relationship Diagram)}{erd_diagram.png}{Vẽ sơ đồ mối quan hệ các bảng hoặc chụp từ SSMS.}

% =========================
% CHƯƠNG 3: PHÂN TÍCH VÀ THIẾT KẾ
% =========================

\chapter{Phân tích và thiết kế}

\section{Khảo sát yêu cầu chi tiết}

\subsection{Yêu cầu chức năng (Functional Requirements)}

\begin{longtable}{|p{0.08\textwidth}|p{0.50\textwidth}|p{0.25\textwidth}|}
\toprule
\textbf{Mã} & \textbf{Mô tả} & \textbf{Phạm vi} \\
\midrule
\endhead

FR-01 & Người dùng public xem danh sách tin tức (phân trang, search) & Public \\
\hline
FR-02 & Người dùng public xem chi tiết tin tức & Public \\
\hline
FR-03 & Người dùng public xem danh sách/chi tiết điểm đến & Public \\
\hline
FR-04 & Người dùng public xem danh sách/chi tiết ẩm thực & Public \\
\hline
FR-05 & Người dùng public gửi feedback & Public \\
\hline
FR-06 & Xem bản đồ hành chính trên trang chủ & Public \\
\hline
FR-07 & Đăng nhập hệ thống admin & Admin/Editor \\
\hline
FR-08 & Quên mật khẩu và đặt lại (self-service) & Admin/Editor \\
\hline
FR-09 & Tạo, chỉnh sửa, xoá tin tức & Admin/Editor \\
\hline
FR-10 & Tạo, chỉnh sửa, xoá điểm đến & Admin/Editor \\
\hline
FR-11 & Tạo, chỉnh sửa, xoá ẩm thực & Admin/Editor \\
\hline
FR-12 & Upload, xoá hình ảnh/media & Admin/Editor \\
\hline
FR-13 & Duyệt và phê duyệt bình luận & Admin \\
\hline
FR-14 & Quản lý cài đặt hệ thống (logo, slogan, map) & Admin \\
\hline
FR-15 & Quản lý người dùng (tạo, xoá, đổi role) & Admin \\
\hline
FR-16 & Admin reset mật khẩu trực tiếp cho user & Admin \\
\hline
FR-17 & Xem log hành động admin & Admin \\
\hline

\bottomrule
\end{longtable}

\subsection{Yêu cầu phi chức năng (Non-Functional Requirements)}

\begin{itemize}[itemsep=3pt]
  \item \textbf{Bảo mật:}
  \begin{itemize}
    \item Tất cả API admin bắt buộc xác thực JWT.
    \item Mật khẩu mã hóa bcrypt.
    \item CORS cấu hình chặt chẽ.
  \end{itemize}
  
  \item \textbf{Hiệu năng:}
  \begin{itemize}
    \item Danh sách dùng phân trang (mặc định 10-20 item).
    \item Media upload tối đa 50MB (tuỳ cấu hình).
    \item Response time < 2 giây cho API public.
  \end{itemize}
  
  \item \textbf{Tính ổn định:}
  \begin{itemize}
    \item Dữ liệu bền vững (persistent) trên Azure SQL.
    \item Tự động migration schema khi update.
    \item Retry mechanism cho các lỗi tạm thời.
  \end{itemize}
  
  \item \textbf{Khả năng triển khai:}
  \begin{itemize}
    \item CI/CD tự động qua GitHub Actions.
    \item Deploy frontend chỉ cần push code.
    \item Deploy backend bằng zip deployment đơn giản.
  \end{itemize}
\end{itemize}

\section{Use-Case Diagram}

\ImagePlaceholder{Use-Case Diagram tổng quan}{usecase_diagram.png}{Vẽ biểu đồ use-case: Public xem nội dung, Editor CRUD, Admin users settings}

\section{Thiết kế kiến trúc hệ thống}

\subsection{Sơ đồ kiến trúc 3-tier}

Hệ thống theo kiến trúc 3 tầng:

\begin{enumerate}[itemsep=3pt]
  \item \textbf{Presentation Layer (Frontend):}
  \begin{itemize}
    \item React SPA chạy trong trình duyệt.
    \item Gọi API backend qua Axios + interceptor.
    \item Quản lý state toàn cục (nếu cần dùng Context API hoặc Redux).
  \end{itemize}
  
  \item \textbf{API/Business Logic Layer (Backend):}
  \begin{itemize}
    \item ASP.NET Core Controllers xử lý request.
    \item Services chứa business logic.
    \item Middleware xử lý authentication, error handling.
  \end{itemize}
  
  \item \textbf{Data Access Layer:}
  \begin{itemize}
    \item Entity Framework Core ORM.
    \item Repositories pattern (tuỳ chọn).
    \item Azure SQL Database backend.
  \end{itemize}
\end{enumerate}

\ImagePlaceholder{Sơ đồ kiến trúc 3-tier}{architecture_3tier.png}{Vẽ hoặc chụp sơ đồ: Frontend SWA API App Service DB Azure SQL}

\subsection{Cấu hình Static Web Apps cho SPA}

Azure Static Web Apps yêu cầu cấu hình rewrite để hỗ trợ SPA routing:

\begin{lstlisting}[language=json,caption={Cấu hình SPA rewrite trong staticwebapp.config.json}]
{
  "navigationFallback": {
    "rewrite": "/index.html",
    "exclude": ["/assets/*", "/api/*"]
  }
}
\end{lstlisting}

Khi người dùng reload page \texttt{/admin/users}, server trả về \texttt{index.html}, đặt để React Router xử lý routing.

\section{Thiết kế API}

\subsection{Base URL và cấu trúc endpoint}

\begin{center}
\begin{tabular}{|l|l|l|}
\hline
\textbf{Nhóm} & \textbf{Base} & \textbf{Ví dụ endpoint} \\
\hline
Public & \texttt{/api/public/} & \texttt{GET /api/news} \\
\hline
Admin & \texttt{/api/admin/} & \texttt{POST /api/admin/news} \\
\hline
Auth & \texttt{/api/auth/} & \texttt{POST /api/auth/login} \\
\hline
Settings & \texttt{/api/settings/} & \texttt{GET /api/settings} \\
\hline
\end{tabular}
\end{center}

\subsection{Endpoint công khai (Public API)}

\begin{lstlisting}[language=bash,caption={Danh sách endpoint public}]
GET  /api/news?page=1&pageSize=10&q=search             # Lấy tin tức
GET  /api/news/{id}                                      # Chi tiết tin tức
GET  /api/news/slug/{slug}                               # Tìm theo slug
GET  /api/destinations?page=1&pageSize=10                # Danh sách địa điểm
GET  /api/destinations/{id}                              # Chi tiết địa điểm
GET  /api/cuisines?page=1&pageSize=10                    # Danh sách ẩm thực
GET  /api/settings                                       # Cài đặt hệ thống
POST /api/feedback                                       # Gửi feedback
\end{lstlisting}

\subsection{Endpoint admin (yêu cầu Bearer token + role Admin/Editor)}

\begin{lstlisting}[language=bash,caption={Danh sách endpoint admin}]
POST   /api/admin/news                                   # Tạo tin tức
GET    /api/admin/news?page=1&pageSize=20&status=1      # Danh sách admin
PUT    /api/admin/news/{id}                              # Cập nhật tin tức
DELETE /api/admin/news/{id}                              # Xoá tin tức

POST   /api/admin/destinations                           # Tạo địa điểm
PUT    /api/admin/destinations/{id}                      # Cập nhật
DELETE /api/admin/destinations/{id}                      # Xoá

POST   /api/admin/media                                  # Upload media
GET    /api/admin/media/{id}                             # Chi tiết file
DELETE /api/admin/media/{id}                             # Xoá file

GET    /api/admin/users                                  # Danh sách user
POST   /api/admin/users                                  # Tạo user
PATCH  /api/admin/users/{id}/role                        # Đổi role
DELETE /api/admin/users/{id}                             # Xoá user
POST   /api/admin/users/{id}/reset-password              # Reset mật khẩu

GET    /api/settings                                     # Lấy settings
PUT    /api/settings/{key}                               # Cập nhật setting

GET    /api/admin/audit-logs                             # Xem log hành động
\end{lstlisting}

\subsection{Endpoint xác thực (Auth API)}

\begin{lstlisting}[language=bash,caption={Endpoint auth}]
POST /api/auth/login                                     # Đăng nhập
POST /api/auth/refresh-token                             # Cấp token mới
POST /api/auth/forgot-password                           # Quên mật khẩu
POST /api/auth/reset-password                            # Đặt lại mật khẩu
\end{lstlisting}

\subsection{Response Format}

Tất cả API trả về JSON với cấu trúc chuẩn:

\begin{lstlisting}[language=json,caption={Response thành công với danh sách}]
{
  "items": [
    {"id": 1, "title": "Tin tức 1", "content": "..."},
    {"id": 2, "title": "Tin tức 2", "content": "..."}
  ],
  "page": 1,
  "pageSize": 10,
  "total": 25,
  "totalPages": 3
}
\end{lstlisting}

\begin{lstlisting}[language=json,caption={Response lỗi}]
{
  "type": "https://tools.ietf.org/html/rfc9110#section-15.5.1",
  "title": "Bad Request",
  "status": 400,
  "detail": "Title is required",
  "traceId": "..."
}
\end{lstlisting}

\ImagePlaceholder{Swagger UI - Tài liệu API interactive}{swagger_ui.png}{Chụp Swagger UI trên máy production, hiển thị 3-4 endpoint chính.}

\section{Thiết kế giao diện (UI/UX)}

\subsection{Public Website}

Trang công khai bao gồm:
\begin{itemize}[itemsep=2pt]
  \item Trang chủ: banner, các tin/điểm đến nổi bật, bản đồ hành chính.
  \item Trang tin tức: danh sách phân trang, search, filter theo ngày.
  \item Trang chi tiết tin: toàn bộ nội dung, comments, share.
  \item Trang điểm đến: map tương tác, chi tiết địa điểm.
  \item Trang liên hệ / feedback form.
\end{itemize}

\subsection{Admin Portal}

Bao gồm:
\begin{itemize}[itemsep=2pt]
  \item Dashboard: thống kê số tin, số user, hoạt động gần đây.
  \item Quản lý tin tức: bảng CRUD với search, filter theo status.
  \item Quản lý địa điểm: CRUD, upload hình.
  \item Quản lý ẩm thực: tương tự.
  \item Quản lý media: danh sách file, xoá, preview.
  \item Quản lý người dùng: bảng user, tạo, đổi role, reset password.
  \item Cài đặt hệ thống: form nhập key-value.
  \item Log hành động: bảng audit log chỉ xem.
\end{itemize}

\ImagePlaceholder{Trang admin dashboard}{ui_admin_dashboard.png}{Chụp trang admin hiển thị thống kê, menu left sidebar}
\ImagePlaceholder{Trang quản lý tin tức Admin}{ui_admin_news.png}{Chụp bảng danh sách tin, nút thêm sửa xoá, filter status}
\ImagePlaceholder{Modal tạo sửa tin tức}{ui_modal_news_edit.png}{Chụp form: title, slug, content editor WYSIWYG status featured image}

% =========================
% CHƯƠNG 4: CÀI ĐẶT/TRIỂN KHAI
% =========================

\chapter{Cài đặt và triển khai}

\section{Cấu trúc mã nguồn}

\subsection{Frontend folder structure}

\begin{lstlisting}[language=bash,caption={Cấu trúc src/ frontend}]
src/
├── app/
│   ├── layouts/           # Layout component
│   ├── pages/             # Trang public + admin
│   │   ├── public/        # Home, News, Destinations, ...
│   │   ├── admin/         # Dashboard, AdminNews, AdminUsers, ...
│   │   ├── auth/          # Login, ResetPassword, ...
│   ├── components/        # Component tái sử dụng
│   ├── services/          # API wrapper (*.service.ts)
│   ├── routes.tsx         # React Router config
│   ├── App.tsx
│   └── main.tsx
├── lib/
│   ├── api.ts             # Axios instance + interceptor
│   └── utils.ts
├── styles/                # Global CSS
└── types/                 # TypeScript interfaces
\end{lstlisting}

\subsection{Backend folder structure}

\begin{lstlisting}[language=bash,caption={Cấu trúc Sadec.Api/ backend}]
Sadec.Api/
├── Controllers/           # API endpoints
│   ├── AdminNewsController.cs
│   ├── AdminUsersController.cs
│   ├── AuthController.cs
│   ├── DestinationsController.cs
│   ├── NewsController.cs
│   └── ...
├── Dtos/                  # Request/Response DTO
│   ├── AuthDtos.cs
│   ├── NewsDtos.cs
│   ├── UserAdminDtos.cs
│   └── ...
├── Entities/              # EF Core entity models
│   ├── ApplicationUser.cs
│   ├── News.cs
│   ├── Destination.cs
│   └── ...
├── Data/
│   ├── ApplicationDbContext.cs
│   ├── DataSeeder.cs
│   └── Migrations/
├── Services/              # Business logic
│   ├── Email/
│   │   ├── IEmailSender.cs
│   │   └── SmtpEmailSender.cs
│   ├── AuditService.cs
│   └── ...
├── Middlewares/           # Custom middleware
│   ├── ErrorHandlingMiddleware.cs
│   └── ...
├── Program.cs             # Startup config
├── appsettings.json       # Configuration
└── Sadec.Api.csproj
\end{lstlisting}

\section{Setup môi trường local}

\subsection{Backend setup}

\begin{lstlisting}[language=bash,caption={Chạy backend local}]
# Clone repo
cd Sadec.Api

# Restore NuGet packages
dotnet restore

# Apply migrations (nếu DB chưa tạo)
dotnet ef database update

# Chạy server
dotnet run --project ./Sadec.Api.csproj

# Swagger: http://localhost:5090/swagger
\end{lstlisting}

\subsection{Frontend setup}

\begin{lstlisting}[language=bash,caption={Chạy frontend local}]
# Cài dependencies
npm ci

# Chạy dev server
npm run dev

# Browser: http://localhost:5173
\end{lstlisting}

\section{Cài đặt reset mật khẩu}

Hệ thống hỗ trợ 2 luồng reset mật khẩu:

\subsection{Luồng 1: Người dùng tự đặt lại (Forgot/Reset Password)}

\subsubsection{Backend implementation}

Endpoint `/api/auth/forgot-password`:

\begin{lstlisting}[language=java,caption={Forgot Password Endpoint}]
[HttpPost]
[AllowAnonymous]
public async Task<IActionResult> ForgotPassword(
    [FromBody] ForgotPasswordRequestDto request)
{
    var user = await _userManager.FindByNameAsync(request.UsernameOrEmail)
        ?? await _userManager.FindByEmailAsync(request.UsernameOrEmail);
    
    if (user == null)
    {
        // Safety: không tiết lộ rằng user không tồn tại
        return Ok(new { 
            message = "If account exists, reset link sent to email" 
        });
    }
    
    // Generate reset token
    var resetToken = await _userManager.GeneratePasswordResetTokenAsync(user);
    
    // Build reset URL
    var resetUrl = $"{Request.Scheme}://{Request.Host}/reset-password?email={user.Email}&token={HttpUtility.UrlEncode(resetToken)}";
    
    // Gửi email (nếu SMTP cấu hình)
    if (_emailSender != null)
    {
        await _emailSender.SendPasswordResetEmailAsync(user.Email, resetUrl);
    }
    
    return Ok(new { 
        message = "Reset link sent. Check email." 
    });
}
\end{lstlisting}

Endpoint `/api/auth/reset-password`:

\begin{lstlisting}[language=java,caption={Reset Password Endpoint}]
[HttpPost]
[AllowAnonymous]
public async Task<IActionResult> ResetPassword(
    [FromBody] ResetPasswordRequestDto request)
{
    var user = await _userManager.FindByEmailAsync(request.Email);
    if (user == null)
        return BadRequest(new { message = "User not found" });
    
    var result = await _userManager.ResetPasswordAsync(
        user, 
        request.Token, 
        request.NewPassword);
    
    if (result.Succeeded)
        return Ok(new { message = "Password reset successful" });
    
    return BadRequest(new { message = string.Join(", ", result.Errors) });
}
\end{lstlisting}

\subsubsection{Frontend implementation}

Trang `/reset-password`:

\begin{lstlisting}[language=bash,caption={URL format}]
/reset-password?email=user@sadec.local&token=ABC123XYZ
\end{lstlisting}

Component xử lý:

\begin{lstlisting}[language=bash,caption={Logic}]
1. Parse URL query params (email, token)
2. Và form nhập mật khẩu mới + xác nhận
3. Gọi POST /api/auth/reset-password { email, token, newPassword }
4. Nếu thành công, chuyển hướng tới /login
\end{lstlisting}

\subsection{Luồng 2: Admin reset mật khẩu trực tiếp}

Endpoint `/api/admin/users/{id}/reset-password`:

\begin{lstlisting}[language=java,caption={Admin Reset Password}]
[HttpPost("{id}/reset-password")]
[Authorize(Roles = "Admin")]
public async Task<IActionResult> ResetPassword(
    string id, 
    [FromBody] UserResetPasswordDto request)
{
    var user = await _userManager.FindByIdAsync(id);
    if (user == null)
        return NotFound();
    
    // Generate token mới
    var resetToken = await _userManager.GeneratePasswordResetTokenAsync(user);
    
    // Reset password
    var result = await _userManager.ResetPasswordAsync(
        user, 
        resetToken, 
        request.NewPassword);
    
    if (result.Succeeded)
    {
        // Log hành động
        _auditService.LogAction("user.password_reset_by_admin", user.Id);
        return NoContent(); // 204
    }
    
    return BadRequest(new { message = string.Join(", ", result.Errors) });
}
\end{lstlisting}

Frontend component (Modal trong Admin Users page):

\begin{lstlisting}[language=bash,caption={UI Modal reset}]
1. Chọn user từ bảng
2. Nhấn "Reset mật khẩu"
3. Modal yêu cầu: mật khẩu mới + xác nhận
4. Validation: độ dài, complexity
5. Gọi API + hiển thị kết quả
\end{lstlisting}

\section{Triển khai trên Azure}

\subsection{Deploy Frontend}

\begin{enumerate}[itemsep=2pt]
  \item Build Vite: \texttt{npm run build} → output \texttt{dist/}
  \item Copy \texttt{staticwebapp.config.json} vào \texttt{dist/}
  \item Push code tới GitHub branch main
  \item GitHub Actions workflow tự động deploy SWA
  \item Xác minh: \url{\FrontendUrl}
\end{enumerate}

\subsection{Deploy Backend}

\begin{enumerate}[itemsep=2pt]
  \item Build Release: \texttt{dotnet publish -c Release -o ./publish\_backend}
  \item Zip folder: \texttt{publish\_backend} → \texttt{publish.zip}
  \item Deploy via Azure CLI:
  \begin{lstlisting}[language=bash,caption={Deploy backend}]
az webapp deploy --resource-group <group> \
  --name <app-name> \
  --src-path ./publish.zip \
  --type zip
  \end{lstlisting}
  \item Xác minh: \url{\SwaggerUrl}
\end{enumerate}

\subsection{Cấu hình App Service}

Các biến môi trường quan trọng (lưu trong App Service Configuration):

\begin{center}
\begin{tabular}{|l|p{5cm}|}
\hline
\textbf{Key} & \textbf{Value ví dụ} \\
\hline
DB\_CONNECTION\_STRING & \texttt{Server=sql.database.windows.net;...} \\
\hline
JWT\_SECRET & \texttt{your-secret-key-min-32-chars} \\
\hline
CORS\_ALLOWED\_ORIGINS & \texttt{https://ashy-field-09c183d00.2.azurestaticapps.net} \\
\hline
SMTP\_HOST & \texttt{smtp.gmail.com} (optional) \\
\hline
SMTP\_PORT & \texttt{587} (optional) \\
\hline
SMTP\_USER & \texttt{your-email@gmail.com} (optional) \\
\hline
FRONTEND\_BASE\_URL & \texttt{https://ashy-field-09c183d00.2.azurestaticapps.net} \\
\hline
\end{tabular}
\end{center}

% =========================
% CHƯƠNG 5: KIỂM THỬ
% =========================

\chapter{Kiểm thử}

\section{Mục tiêu kiểm thử}

\begin{itemize}[itemsep=2pt]
  \item Xác nhận các luồng chính hoạt động đúng.
  \item Xác nhận bảo mật: API admin bắt buộc xác thực, CORS hoạt động.
  \item Xác nhận dữ liệu bền vững: không bị reset khi restart.
  \item Xác nhận SPA routing không lỗi 404.
\end{itemize}

\section{Kiểm thử thủ công}

\subsection{Test Cases}

\begin{longtable}{|p{0.06\textwidth}|p{0.45\textwidth}|p{0.22\textwidth}|p{0.20\textwidth}|}
\toprule
\textbf{TC} & \textbf{Bước kiểm thử} & \textbf{Kỳ vọng} & \textbf{Kết quả} \\
\midrule
\endhead

TC01 & Truy cập trang chủ public & 200 OK, hiển thị nội dung & PASS \\
\hline

TC02 & Search tin tức & Danh sách tin phù hợp & PASS \\
\hline

TC03 & Đăng nhập admin với credential đúng & Chuyển hướng /admin, nhận token & PASS \\
\hline

TC04 & Đăng nhập admin với sai password & 401 Unauthorized & PASS \\
\hline

TC05 & Gọi API admin mà không token & 401 Unauthorized & PASS \\
\hline

TC06 & Gọi API admin với token nhưng role không đủ & 403 Forbidden & PASS \\
\hline

TC07 & Refresh token khi hết hạn & Tự động lấy token mới, retry thành công & PASS \\
\hline

TC08 & Reload trang /admin/users & Không lỗi 404 (SPA routing hoạt động) & PASS \\
\hline

TC09 & Tạo tin tức mới (admin) & Tin được lưu, xuất hiện trong danh sách admin & PASS \\
\hline

TC10 & Xoá tin tức (admin) & Tin bị xoá, không còn hiển thị & PASS \\
\hline

TC11 & Upload media (admin) & File lưu, URL trả về, download OK & PASS \\
\hline

TC12 & Admin tạo user mới & User được tạo, có email + sinh password tạm & PASS \\
\hline

TC13 & Admin reset password user & User đăng nhập được với password mới & PASS \\
\hline

TC14 & Quên mật khẩu (không SMTP) & Endpoint trả 200 (tuy không gửi email) & PASS \\
\hline

TC15 & Reset password bằng token & Mật khẩu update, đăng nhập với PW mới OK & PASS \\
\hline

TC16 & Restart app, dữ liệu đã tạo vẫn có & Dữ liệu bền vững trên Azure SQL & PASS \\
\hline

\bottomrule
\end{longtable}

\section{Kiểm thử API bằng Swagger/Postman}

Truy cập \url{\SwaggerUrl} để test interactive:

\begin{itemize}[itemsep=2pt]
  \item Login endpoint: POST /api/auth/login
  \item Copy token từ response
  \item Click "Authorize" trong Swagger, paste token
  \item Test các endpoint admin (News, Users, Settings)
\end{itemize}

\ImagePlaceholder{Swagger Test Login endpoint}{test_swagger_login.png}{Chụp Swagger request body response với token}
\ImagePlaceholder{Swagger Test Admin Reset Password}{test_swagger_admin_reset.png}{Chụp Swagger POST api admin users id reset-password trả 204}

\section{Kiểm thử Giao diện}

\subsection{Kiểm thử Login}

\begin{itemize}[itemsep=2pt]
  \item Input credential admin@sadec.local / Admin@123
  \item Xác nhận redirect /admin
  \item Xác nhận token lưu localStorage
\end{itemize}

\subsection{Kiểm thử Reset Password}

\begin{itemize}[itemsep=2pt]
  \item Trang /reset-password?email=user@sadec.local&token=XYZ
  \item Nhập password mới 2 lần
  \item Submit → xác nhận redirect /login
  \item Đăng nhập lại với password mới
\end{itemize}

\subsection{Kiểm thử Admin Users}

\begin{itemize}[itemsep=2pt]
  \item Xem danh sách user (bảng phân trang)
  \item Nhấn "Thêm người dùng": modal mở
  \item Nhập email + password, nhấn "Tạo"
  \item Xác nhận user mới xuất hiện trong danh sách
  \item Nhấn "Reset mật khẩu" trên user: modal mở
  \item Nhập password mới, nhấn "Xác nhận"
  \item Xác nhận thành công / lỗi hiển thị rõ
\end{itemize}

\ImagePlaceholder{UI form Thêm người dùng Admin}{ui_form_add_user.png}{Chụp modal email field password field nút Tạo}
\ImagePlaceholder{UI modal Reset mật khẩu Admin Users}{ui_modal_reset_password.png}{Chụp modal hiển thị user ô password mới ô xác nhận nút Xác nhận}

\section{Vấn đề phát hiện và xử lý}

\subsection{Vấn đề 1: SPA reload trang con bị 404}

\textbf{Nguyên nhân:} Server không cấu hình rewrite cho SPA.

\textbf{Giải pháp:} Thêm \texttt{navigationFallback} trong \texttt{staticwebapp.config.json}.

\textbf{Xác minh:} Reload trang /admin/users → không 404, hiển thị bình thường.

\subsection{Vấn đề 2: Dữ liệu bị reset sau restart}

\textbf{Nguyên nhân:} Dùng InMemory Database.

\textbf{Giải pháp:} Chuyển sang Azure SQL + EF Core migrations.

\textbf{Xác minh:} Restart app, dữ liệu vẫn có.

\subsection{Vấn đề 3: Admin add-user không rõ lỗi}

\textbf{Nguyên nhân:} Frontend không đọc backend error message.

\textbf{Giải pháp:} Đổi service dùng shared API client, hiển thị \texttt{response.data.message}.

\textbf{Xác minh:} Tạo user với email trùng → hiển thị "Email already exists".

\subsection{Vấn đề 4: Quên mật khẩu không gửi Gmail}

\textbf{Nguyên nhân:} SMTP chưa cấu hình.

\textbf{Giải pháp:} Bổ sung admin reset password; SMTP là tuỳ chọn.

\textbf{Xác minh:} Endpoint trả 200; nếu có SMTP sẽ gửi email.

% =========================
% CHƯƠNG 6: KẾT LUẬN
% =========================

\chapter{Kết luận}

\section{Tóm tắt kết quả}

Dự án \ProjectName{} đã hoàn thành thành công với các thành tựu chính:

\begin{itemize}[itemsep=3pt]
  \item \textbf{Hệ thống hoàn chỉnh:} Website public + admin portal với đầy đủ CRUD nội dung.
  
  \item \textbf{Bảo mật:} Xác thực JWT, phân quyền RBAC, CORS cấu hình chặt.
  
  \item \textbf{Tính bền vững:} Dữ liệu lưu Azure SQL, migrations tự động, không bị reset.
  
  \item \textbf{Triển khai thực tế:} Frontend trên SWA, backend trên App Service, CI/CD qua GitHub Actions.
  
  \item \textbf{Tính năng nâng cao:} Forgot/reset password, admin reset password, audit logging, settings management.
  
  \item \textbf{Quy trình phát triển:} Cấu trúc code rõ ràng, dễ bảo trì và mở rộng.
\end{itemize}

\section{Kỹ năng và công nghệ học được}

\begin{itemize}[itemsep=2pt]
  \item \textbf{Frontend:} React, TypeScript, Vite, React Router, Axios interceptor.
  \item \textbf{Backend:} ASP.NET Core, EF Core, ASP.NET Identity, JWT authentication.
  \item \textbf{Cloud:} Azure Static Web Apps, Azure App Service, Azure SQL, deployment strategies.
  \item \textbf{DevOps:} GitHub Actions CI/CD, zip deployment, environment configuration.
\end{itemize}

\section{Hạn chế}

\begin{itemize}[itemsep=2pt]
  \item SMTP email delivery chưa cấu hình (optional feature).
  \item Unit tests/integration tests còn hạn chế (có thể bổ sung thêm).
  \item Monitoring/logging advanced (Application Insights) chưa áp dụng.
  \item Performance optimization (caching, CDN) là hướng phát triển tiếp theo.
\end{itemize}

\section{Hướng phát triển tương lai}

\begin{enumerate}[itemsep=3pt]
  \item \textbf{Hoàn thiện email:} Cấu hình SMTP + email templates chuyên nghiệp.
  
  \item \textbf{Monitoring:} Tích hợp Application Insights, alert khi error.
  
  \item \textbf{Caching:} Redis cache cho những query thường xuyên.
  
  \item \textbf{Performance:} CDN cho static assets, database indexing tuning.
  
  \item \textbf{Tính năng:} Two-factor authentication, OAuth social login.
  
  \item \textbf{Testing:} Unit tests, integration tests, E2E tests.
  
  \item \textbf{SEO:} Metadata, sitemap, structured data for content.
\end{enumerate}

\section{Lời kết}

Đề tài này là một ứng dụng web thực tế, đã chuẩn bị theo quy trình DevOps chuyên nghiệp.
Quá trình phát triển từ thiết kế, cài đặt, kiểm thử đến triển khai cloud đã minh chứng rõ
ứng dụng của các kiến thức về web development, cloud computing, và software engineering.

Hy vọng báo cáo này cung cấp nền tảng vững chắc cho các nhân viên, teacher và các bên liên quan
để hiểu rõ kiến trúc, công nghệ, và quy trình phát triển của hệ thống.

% =========================
% Tài liệu tham khảo
% =========================

\begin{thebibliography}{9}

\bibitem{aspnet}
Microsoft, \textit{ASP.NET Core Documentation}. 
\url{https://learn.microsoft.com/aspnet/core/}

\bibitem{identity}
Microsoft, \textit{ASP.NET Core Identity Documentation}. 
\url{https://learn.microsoft.com/aspnet/core/security/authentication/identity}

\bibitem{efcore}
Microsoft, \textit{Entity Framework Core}. 
\url{https://learn.microsoft.com/ef/core/}

\bibitem{react}
Facebook Developers, \textit{React Official Documentation}. 
\url{https://react.dev/}

\bibitem{swa}
Microsoft, \textit{Azure Static Web Apps}. 
\url{https://learn.microsoft.com/azure/static-web-apps/}

\bibitem{appservice}
Microsoft, \textit{Azure App Service Documentation}. 
\url{https://learn.microsoft.com/azure/app-service/}

\bibitem{azuresql}
Microsoft, \textit{Azure SQL Database Documentation}. 
\url{https://learn.microsoft.com/azure/azure-sql/}

\bibitem{jwt}
Auth0, \textit{JWT.io - Introduction to JSON Web Tokens}. 
\url{https://jwt.io/introduction}

\bibitem{rest}
Roy Fielding, \textit{Representational State Transfer (REST)}. 
\url{https://www.ics.uci.edu/~fielding/pubs/dissertation/rest_arch_style.htm}

\end{thebibliography}

% =========================
% PHỤ LỤC
% =========================

\appendix

\chapter{Phụ lục A: Danh sách ảnh cần chèn}

\begin{longtable}{|p{0.30\textwidth}|p{0.65\textwidth}|}
\toprule
\textbf{Tên ảnh (figures/)} & \textbf{Nội dung gợi ý} \\
\midrule
ui_homepage.png & Trang chủ public: banner, tin nổi bật, bản đồ hành chính. \\
\hline
erd_diagram.png & Entity Relationship Diagram các bảng (SSMS hoặc draw.io). \\
\hline
usecase_diagram.png & Use-case diagram: Public actor, Editor, Admin. \\
\hline
architecture_3tier.png & Sơ đồ 3-tier: Frontend (SWA) -> API (App Service) -> DB. \\
\hline
swagger_ui.png & Swagger UI trên production: endpoint chính. \\
\hline
ui_admin_dashboard.png & Trang admin dashboard: menu, thống kê. \\
\hline
ui_admin_news.png & Bảng danh sách tin tức admin: header, data rows, action buttons. \\
\hline
ui_modal_news_edit.png & Modal tạo/sửa tin: form fields, buttons. \\
\hline
ui_form_add_user.png & Form thêm người dùng: email, password fields. \\
\hline
ui_modal_reset_password.png & Modal reset mật khẩu: user info, password fields, button. \\
\hline
test_swagger_login.png & Swagger: POST /api/auth/login + response. \\
\hline
test_swagger_admin_reset.png & Swagger: Endpoint reset password admin. \\
\hline
\end{longtable}

\chapter{Phụ lục B: Code Snippets Chiều dài}

\section{Frontend - API Client Setup}

\begin{lstlisting}[language=js,caption={src/lib/api.ts - Axios Interceptor}]
import axios, { AxiosInstance } from "axios";

const API_BASE_URL = 
  import.meta.env.VITE_API_URL || 
  (window.location.hostname.includes('azurestaticapps.net') 
    ? 'https://sadec-backend.azurewebsites.net/api' 
    : 'http://localhost:5000/api');

export const api = axios.create({
  baseURL: API_BASE_URL,
});

// Request interceptor: attach token
api.interceptors.request.use(config => {
  const token = localStorage.getItem("token");
  if (token) {
    config.headers.Authorization = `Bearer ${token}`;
  }
  return config;
});

// Response interceptor: refresh token on 401
api.interceptors.response.use(
  response => response,
  async error => {
    const originalRequest = error.config;
    if (error.response?.status === 401 && !originalRequest._retry) {
      originalRequest._retry = true;
      const refreshToken = localStorage.getItem("refreshToken");
      if (refreshToken) {
        const resp = await axios.post(`${API_BASE_URL}/auth/refresh-token`, {
          token: localStorage.getItem("token"),
          refreshToken
        });
        localStorage.setItem("token", resp.data.token);
        originalRequest.headers.Authorization = `Bearer ${resp.data.token}`;
        return api(originalRequest);
      }
    }
    return Promise.reject(error);
  }
);
\end{lstlisting}

\section{Backend - Authentication Middleware}

\begin{lstlisting}[language=java,caption={Program.cs - JWT Setup}]
var builder = WebApplicationBuilder.CreateBuilder(args);

// Add Authentication
builder.Services.AddAuthentication(JwtBearerDefaults.AuthenticationScheme)
.AddJwtBearer(opts => {
  opts.TokenValidationParameters = new()
  {
    ValidateIssuer = false,
    ValidateAudience = false,
    ValidateLifetime = true,
    ValidateIssuerSigningKey = true,
    IssuerSigningKey = new SymmetricSecurityKey(
      Encoding.UTF8.GetBytes(builder.Configuration["Jwt:Secret"] ?? ""))
  };
});

// Add Authorization
builder.Services.AddAuthorization();

// Add DbContext
builder.Services.AddDbContext<ApplicationDbContext>(opts =>
  opts.UseSqlServer(builder.Configuration.GetConnectionString("DefaultConnection")));

var app = builder.Build();

app.UseAuthentication();
app.UseAuthorization();

app.MapControllers();
app.Run();
\end{lstlisting}

\chapter{Phụ lục C: Các lệnh thường dùng}

\section{Frontend}

\begin{lstlisting}[language=bash,caption={Lệnh frontend}]
# Cài dependencies
npm ci

# Chạy dev server
npm run dev

# Build production
npm run build

# Preview build
npm run preview

# Lint (nếu có ESLint)
npm run lint
\end{lstlisting}

\section{Backend}

\begin{lstlisting}[language=bash,caption={Lệnh backend}]
# Restore packages
dotnet restore

# Build
dotnet build

# Chạy dev
dotnet run

# Publish Release
dotnet publish -c Release -o ./publish_backend

# EF Core migrations
dotnet ef migrations add <MigrationName>

# Apply migrations
dotnet ef database update

# View schema
dotnet ef dbcontext scaffold ...
\end{lstlisting}

\chapter{Phụ lục D: Troubleshooting}

\section{Frontend}

\begin{itemize}[itemsep=2pt]
  \item \textbf{CORS error:} Backend cần cấu hình CORS cho origin của frontend.
  \item \textbf{Token không gửi:} Kiểm tra localStorage có token, axios interceptor chạy đúng.
  \item \textbf{404 khi reload route:} Kiểm tra staticwebapp.config.json có navigationFallback.
\end{itemize}

\section{Backend}

\begin{itemize}[itemsep=2pt]
  \item \textbf{Connection string error:} Kiểm tra App Service configuration có DB connection string.
  \item \textbf{Migration fails:} Chắc chắn database đã tạo, firewall rule cho App Service.
  \item \textbf{401 Unauthorized:} Token hết hạn hoặc sai secret key.
\end{itemize}

\end{document}
